% ============================================================
% M3 S3 练习件·W2 臂B —— 课时09 圆与圆的位置关系（2.3.4）＝照抄母版件（课时01）体例
% 生成：M3 成卷轮 S3 W2 臂B｜2026-09-14｜编译 cwd＝本目录（中文路径禁 \input 绝对路径）
%   xelatex main-true.tex ×2 → 含详解印本；xelatex main-false.tex ×2 → 纯题版（单源双本）
% 输入链（全只读）：题面库 课时09-圆与圆的位置关系.md（题面逐字装配）＋ -答案侧.md（值逐字回嵌）
%   ＋ 定稿/批B-课时09-圆与圆的位置关系.md（详解＝件内【亲算】格回嵌源）；toolchain＝qp-m3.sty
%   （钉后版 md5 7c3930362be8a0a2bdf21bbf8ac16573，本地挂载副本，破损即停）。
% 母版体例：详见《练习件母版体例说明.md》（课时01 件）；门谱读数见 过程件 工作区/_tmpM3S3W2臂B0914/。
% 键账：件内 ansblock 键集 ≡ 题面库片练键集 ≡ manifest 练键序 T1~T16（16 键，零缺漏零多余）；
%   拓区隔离：2章-拓/测/滚 键不入本件（S4 拓展册收）；导学键 G1~G5 属导学件域，亦不入。
% 卷式例外案B：练习件不属卷式——答案块物理落点恒为题后 ansblock，禁卷末附卷制；
%   全线括线模（导言 \ansblockgrayfalse 一次置定，件内不切换；灰底盒不可跨栏断，练习件禁用）。
% 题序＝manifest 键序 T1~T16（零重排）；印面题号 1~16 连号（\tihao 号＝\ansitem 号同键）；
%   三档切位 1~7／8~14／15~16（照库序切位，非难度重排）。
% 槽配实录：单选6（T1~T6）＋多选4（T7~T10，＝manifest 多选数，多选门≤4 触顶合规
%   ——题面库片头与 manifest 均已登记「不得再增」）＋填空4（T11~T14）＋解答2（T15/T16）；
%   10选4填2解 为波次方案基线槽配，本片实配 6单选＋4多选＋4填＋2解 照题面侧头标
%   （批B §四 满足度表同口径：源件12 多选富集所致，已在批B 台账注记）。
% 选项槽位：默认四连排（0.25\linewidth−0.25\lxhang），超宽降档梯 四连排→两连排
%   （0.5\linewidth−0.5\lxhang−0.25em）→单列 \lxopt；禁缩字号/负 kern/删标点凑宽。
%   本片实测降档（槽宽门读数 0914）：T1 单列（D槽35.8mm＞31.69）、T2 两连排（C槽19.0mm）、
%   T6 两连排（B槽21.0mm）、T7~T9 单列（\lxopt 长陈述项），余四连排（T3 13.0／T4 11.1／T5 14.9／T10 11.3）。
% 解答书写区：T15/T16 均无小问，\liubai[16mm]（默认档；纯题档吞 ansblock 不吞 \liubai）。
% 填空印答：T11~T14 空位 \kongda{值}（详解版印答、纯题版退定宽空线，两档等形）。
% 难度词：题面侧头标第 4 字段系裸数值（批B 口径），派生标注照同章导学件09 同族口径
%   0.94/0.85→简单、0.65→中档（派生登记见值台账）；知识点 N 照导学件09 知识导学
%   例/变式预排：一＝位置关系的判定与公切线、二＝公共弦·切点弦·定圆与证明、
%   三＝存在性·对称·综合判定（T3/T12 归知识点三，T14 归知识点二，同导学件例位）。
% 印面清洗：T9/T12 题面侧【注】行不入印面；T14 题面「⊙C」照同章导学件09 变式14 同款
%   规范化为「圆 C」（\odot 免路由，批B 题面同式）；行尾句点照导学件口径统一全角。
% ============================================================
\documentclass[fontset=none]{ctexart}
\usepackage{qp-m3}
% —— 印面逐字符号路由补钉（承导学件母版；− 走 TNR、▱ NSC 子块；禁改 sty 保 md5） ——
\xeCJKDeclareCharClass{Default}{"2212}%
\setCJKfamilyfont{nscblk}[Path=C:/提示词/工作区/字替对照-0909/variantF/fonts/,BoldFont=NSC-w600.ttf]{NSC-w500.ttf}%
\xeCJKDeclareSubCJKBlock{nscblk}{"25B1}%
\newcommand{\pxparallelogram}{{\CJKfamily{nscblk}▱}}%
% —— 断行微调（承导学件母版实测档，三零门承重；\emergencystretch 不另设） ——
\xeCJKsetup{CJKglue={\hskip 0.10em plus 0.08em minus 0.05em}}%
% —— 练习件全线括线模（S3 波次方案 §四.3：导言一次置定、件内不切换） ——
\ansblockgrayfalse
% —— 页脚件名（qp-headfoot 复用入口） ——
\renewcommand{\qpzhangming}{第二章\quad 平面解析几何}%
\renewcommand{\qpjianming}{练习件}%

\begin{document}
%装配轮S6三本跨本连续页码（G7方案B）setcounter 注入位
\setcounter{page}{96}

\keshi{第9课时\quad 圆与圆的位置关系}
\begin{multicols}{2}
\emergencystretch=1em

% ============================================================
% 基础巩固（题1~7＝T1~T7：单选6＋多选1；题后紧跟答案制，全线括线 ansblock）
% ============================================================
\dangbadge{基}{础}{巩}{固}

\tihao{1}\zu{圆与圆的位置关系×单选} \tieside{简单(知识点一)}已知圆\(C_1\)：\(x^{2}+(y+m)^{2}=2\)与圆\(C_2\)：\((x-m)^{2}+y^{2}=8\)恰有两条公切线，则实数\(m\)的取值范围是\nobreak(\kongwei)

\lxopt{A．\(1<m<3\)}

\lxopt{B．\(-1<m<1\)}

\lxopt{C．\(m>3\)}

\lxopt{D．\(-3<m<-1\)或\(1<m<3\)}

\begin{ansblock}[2章-练-课时09-T1]
% ans:2章-练-课时09-T1
\ansitem{1}{D}
\ansnote{详解}{两条公切线\(\Longleftrightarrow\)两圆相交．圆心\((0,-m)\)，\((m,0)\)，半径\(\sqrt{2}\)，\(2\sqrt{2}\)，圆心距\(d=\sqrt{2}|m|\)．相交\(\Longleftrightarrow 2\sqrt{2}-\sqrt{2}<\sqrt{2}|m|<2\sqrt{2}+\sqrt{2}\)，即\(1<|m|<3\)．故选：D．}
\end{ansblock}

\tihao{2}\zu{圆与圆的位置关系×单选} \tieside{简单(知识点一)}集合\(M=\{(x,y)\mid x^{2}+y^{2}\leqslant 4\}\)，\(N=\{(x,y)\mid(x-1)^{2}+(y-1)^{2}\leqslant r^{2},\ r>0\}\)，且\(M\cap N=N\)，则\(r\)的取值范围是\nobreak(\kongwei)

\bindopt {\fontsize{10.5pt}{\qpTJgLeadLiti}\selectfont \makebox[\dimexpr0.5\linewidth-0.5\lxhang-0.25em\relax][l]{A．\((0,\sqrt{2}-1)\)}\makebox[\dimexpr0.5\linewidth-0.5\lxhang-0.25em\relax][l]{B．\((0,1]\)}\par}
\bindopt {\fontsize{10.5pt}{\qpTJgLeadLiti}\selectfont \makebox[\dimexpr0.5\linewidth-0.5\lxhang-0.25em\relax][l]{C．\((0,2-\sqrt{2}]\)}\makebox[\dimexpr0.5\linewidth-0.5\lxhang-0.25em\relax][l]{D．\((0,2]\)}\par}

\begin{ansblock}[2章-练-课时09-T2]
% ans:2章-练-课时09-T2
\ansitem{2}{C}
\ansnote{详解}{\(M\cap N=N\Longleftrightarrow N\subseteq M\)，即圆\(N\)内含或内切于圆\(M\)：\(d+r\leqslant 2\)．\(d=\sqrt{2}\)，故\(\sqrt{2}+r\leqslant 2\)，得\(0<r\leqslant 2-\sqrt{2}\)．故选：C．}
\end{ansblock}

\tihao{3}\zu{圆与圆的位置关系×单选} \tieside{中档(知识点三)}已知点\(A(-2,0)\)，\(B(2,0)\)，若圆\((x-3)^{2}+y^{2}=r^{2}\)（\(r>0\)）上存在点\(P\)（不同于点\(A\)，\(B\)），使得\(PA\cdot PB=0\)，则实数\(r\)的取值范围是\nobreak(\kongwei)
\bindopt {\fontsize{10.5pt}{\qpTJgLeadLiti}\selectfont \makebox[\dimexpr0.25\linewidth-0.25\lxhang\relax][l]{A．\((1,5)\)}\makebox[\dimexpr0.25\linewidth-0.25\lxhang\relax][l]{B．\([1,5]\)}\makebox[\dimexpr0.25\linewidth-0.25\lxhang\relax][l]{C．\((1,3]\)}\makebox[\dimexpr0.25\linewidth-0.25\lxhang\relax][l]{D．\([3,5)\)}\par}

\begin{ansblock}[2章-练-课时09-T3]
% ans:2章-练-课时09-T3
\ansitem{3}{A}
\ansnote{详解}{\(PA\cdot PB=0\Longleftrightarrow P\)在以\(AB\)为直径的圆\(x^{2}+y^{2}=4\)上（除\(A\)、\(B\)）．存在这样的\(P\Longleftrightarrow\)两圆相交：\(|r-2|<3<r+2\)，得\(r>1\)且\(r<5\)．故选：A．（恰为\(A\)、\(B\)的情形对应相切\(r=1\)或\(r=5\)，已由「不同于\(A\)，\(B\)」排除．）}
\end{ansblock}

\tihao{4}\zu{圆与圆的位置关系×单选} \tieside{简单(知识点一)}若点\(A(1,0)\)和点\(B(4,0)\)到直线\(l\)的距离依次为1和2，则这样的直线有\nobreak(\kongwei)
\bindopt {\fontsize{10.5pt}{\qpTJgLeadLiti}\selectfont \makebox[\dimexpr0.25\linewidth-0.25\lxhang\relax][l]{A．0条}\makebox[\dimexpr0.25\linewidth-0.25\lxhang\relax][l]{B．1条}\makebox[\dimexpr0.25\linewidth-0.25\lxhang\relax][l]{C．2条}\makebox[\dimexpr0.25\linewidth-0.25\lxhang\relax][l]{D．3条}\par}

\begin{ansblock}[2章-练-课时09-T4]
% ans:2章-练-课时09-T4
\ansitem{4}{D}
\ansnote{详解}{满足条件的\(l\)即分别以\(A\)、\(B\)为圆心，1、2为半径的两圆的公切线．圆心距\(|AB|=3=1+2\)，两圆外切，公切线有3条（2外＋1内）．故选：D．}
\end{ansblock}

\tihao{5}\zu{两圆公共弦与根轴×单选} \tieside{简单(知识点二)}已知圆\(O_1\)：\(x^{2}+y^{2}=4\)与圆\(O_2\)：\(x^{2}+6x+y^{2}=0\)相交于点\(A\)，\(B\)，则四边形\(AO_1BO_2\)的面积是\nobreak(\kongwei)
\bindopt {\fontsize{10.5pt}{\qpTJgLeadLiti}\selectfont \makebox[\dimexpr0.25\linewidth-0.25\lxhang\relax][l]{A．\(4\sqrt{2}/3\)}\makebox[\dimexpr0.25\linewidth-0.25\lxhang\relax][l]{B．\(2\sqrt{2}\)}\makebox[\dimexpr0.25\linewidth-0.25\lxhang\relax][l]{C．\(4\sqrt{2}\)}\makebox[\dimexpr0.25\linewidth-0.25\lxhang\relax][l]{D．\(8\sqrt{2}/3\)}\par}

\begin{ansblock}[2章-练-课时09-T5]
% ans:2章-练-课时09-T5
\ansitem{5}{C}
\ansnote{详解}{两方程相减得公共弦\(AB\)：\(x=-2/3\)．\(|AB|=2\sqrt{4-4/9}=8\sqrt{2}/3\)，\(|O_1O_2|=3\)．对角线互相垂直的四边形面积\(=(1/2)\times(8\sqrt{2}/3)\times 3=4\sqrt{2}\)．故选：C．}
\end{ansblock}

\tihao{6}\zu{两圆公共弦与根轴×单选} \tieside{简单(知识点二)}过圆\(M\)：\((x-1)^{2}+y^{2}=4\)内一点\(A(2,1)\)作一弦交圆于\(B\)、\(C\)两点，过点\(B\)、\(C\)分别作圆的切线\(PB\)、\(PC\)，两切线交于点\(P\)，则点\(P\)的轨迹方程为\nobreak(\kongwei)

\bindopt {\fontsize{10.5pt}{\qpTJgLeadLiti}\selectfont \makebox[\dimexpr0.5\linewidth-0.5\lxhang-0.25em\relax][l]{A．\(y-5=0\)}\makebox[\dimexpr0.5\linewidth-0.5\lxhang-0.25em\relax][l]{B．\(x+y+5=0\)}\par}
\bindopt {\fontsize{10.5pt}{\qpTJgLeadLiti}\selectfont \makebox[\dimexpr0.5\linewidth-0.5\lxhang-0.25em\relax][l]{C．\(x+y-5=0\)}\makebox[\dimexpr0.5\linewidth-0.5\lxhang-0.25em\relax][l]{D．\(x-y-5=0\)}\par}

\begin{ansblock}[2章-练-课时09-T6]
% ans:2章-练-课时09-T6
\ansitem{6}{C}
\ansnote{详解}{设\(P(x_0,y_0)\)．以\(MP\)为直径的圆过切点\(B\)、\(C\)（\(MB\perp PB\)），其方程为\((x-1)(x-x_0)+y(y-y_0)=0\)，与圆\(M\)：\(x^{2}+y^{2}-2x-3=0\)相减得公共弦\(BC\)所在直线：\par\noindent \((x_0-1)x+y_0y-x_0-3=0\)．\(A(2,1)\)在\(BC\)上：\(2(x_0-1)+y_0-x_0-3=0\)，即\(x_0+y_0=5\)．故\(P\)的轨迹为\(x+y-5=0\)．故选：C．}
\end{ansblock}

\tihao{7}\duoxuan\hspace{0.6em}\zu{圆与圆的位置关系×多选} \tieside{简单(知识点一)}已知圆\(M\)：\((x-2)^{2}+(y-1)^{2}=1\)，圆\(N\)：\((x+2)^{2}+(y+1)^{2}=1\)，则下列是\(M\)，\(N\)两圆公切线的直线方程为\nobreak(\kongwei)

\lxopt{A．\(y=0\)}

\lxopt{B．\(3x-4y=0\)}

\lxopt{C．\(x-2y+\sqrt{5}=0\)}

\lxopt{D．\(x-2y-\sqrt{5}=0\)}

\begin{ansblock}[2章-练-课时09-T7]
% ans:2章-练-课时09-T7
\ansitem{7}{ACD}
\ansnote{详解}{圆心\(M(2,1)\)，\(N(-2,-1)\)，\(d=2\sqrt{5}>2\)，两圆相离，公切线4条．逐项验圆心到直线的距离是否等于1：A：\(y=0\)到\(M\)、\(N\)的距离均为1，是；B：\(3x-4y=0\)到\(M\)的距离\(|6-4|/5=2/5\neq 1\)，不是；C：\(x-2y+\sqrt{5}=0\)到\(M\)、\(N\)的距离均为\(\sqrt{5}/\sqrt{5}=1\)，是；D：同理距离为1，是．故选：ACD．}
\end{ansblock}

% ============================================================
% 综合提升（题8~14＝T8~T14：多选3（T8~T10）＋填空4（T11~T14））
% ============================================================
\dangbadge{综}{合}{提}{升}

\tihao{8}\duoxuan\hspace{0.6em}\zu{圆与圆的位置关系×多选} \tieside{中档(知识点一)}已知\(C\)：\(x^{2}+y^{2}-6x=0\)，则下述正确的是\nobreak(\kongwei)

\lxopt{A．圆\(C\)的半径\(r=3\)}

\lxopt{B．点\((1,2\sqrt{2})\)在圆\(C\)的内部}

\lxopt{C．直线\(l\)：\(x+\sqrt{3}y+3=0\)与圆\(C\)相切}

\lxopt{D．圆\(C'\)：\((x+1)^{2}+y^{2}=4\)与圆\(C\)相交}

\begin{ansblock}[2章-练-课时09-T8]
% ans:2章-练-课时09-T8
\ansitem{8}{ACD}
\ansnote{详解}{圆即\((x-3)^{2}+y^{2}=9\)．A：\(r=3\)，正确；B：点\((1,2\sqrt{2})\)到圆心\((3,0)\)的距离\(\sqrt{4+8}=2\sqrt{3}>3\)，在圆外，错误；C：\(d=|3+0+3|/2=3=r\)，相切，正确；D：圆心距\(|3-(-1)|=4\)，半径3、2，\(1<4<5\)，相交，正确．故选：ACD．}
\end{ansblock}

\tihao{9}\duoxuan\hspace{0.6em}\zu{圆与圆的位置关系×多选} \tieside{简单(知识点一)}已知圆\(O\)：\(x^{2}+y^{2}=4\)和圆\(C\)：\((x-2)^{2}+(y-3)^{2}=1\)．现给出如下结论，其中正确的是\nobreak(\kongwei)

\lxopt{A．圆\(O\)与圆\(C\)有四条公切线}

\lxopt{B．过\(C\)且在两坐标轴上截距相等的直线方程为\(x+y=5\)或\(x-y+1=0\)}

\lxopt{C．过\(C\)且与圆\(O\)相切的直线方程为\(9x-16y+30=0\)}

\lxopt{D．\(P\)、\(Q\)分别为圆\(O\)和圆\(C\)上的动点，则\(|PQ|\)的最大值为\(\sqrt{13}+3\)，最小值为\(\sqrt{13}-3\)}

\begin{ansblock}[2章-练-课时09-T9]
% ans:2章-练-课时09-T9
\ansitem{9}{AD}
\ansnote{详解}{\(|OC|=\sqrt{4+9}=\sqrt{13}>2+1=3\)，两圆外离，公切线4条，A正确．\par\noindent B：\(x+y=5\)过\(C(2,3)\)且两截距均为5；\(x-y+1=0\)过\(C\)，但两截距为\(-1\)与\(1\)，不相等，故B错误．\par\noindent C：\(C\)在圆\(O\)外（\(|OC|=\sqrt{13}>2\)），过\(C\)的切线应有两条，选项只列一条且表述为唯一，不成立，C错误．\par\noindent D：\(|PQ|\)的最大值为\(|OC|+2+1=\sqrt{13}+3\)，最小值为\(|OC|-2-1=\sqrt{13}-3\)，D正确．故选：AD．}
\end{ansblock}

\tihao{10}\duoxuan\hspace{0.6em}\zu{圆与圆的位置关系×多选} \tieside{简单(知识点一)}已知圆\(C_1\)：\(x^{2}+(y-a)^{2}=9\)与圆\(C_2\)：\((x-a)^{2}+y^{2}=1\)有四条公切线，则实数\(a\)的取值可能是\nobreak(\kongwei)
\bindopt {\fontsize{10.5pt}{\qpTJgLeadLiti}\selectfont \makebox[\dimexpr0.25\linewidth-0.25\lxhang\relax][l]{A．\(-4\)}\makebox[\dimexpr0.25\linewidth-0.25\lxhang\relax][l]{B．\(-2\)}\makebox[\dimexpr0.25\linewidth-0.25\lxhang\relax][l]{C．\(2\sqrt{2}\)}\makebox[\dimexpr0.25\linewidth-0.25\lxhang\relax][l]{D．\(3\)}\par}

\begin{ansblock}[2章-练-课时09-T10]
% ans:2章-练-课时09-T10
\ansitem{10}{AD}
\ansnote{详解}{四条公切线\(\Longleftrightarrow\)两圆外离．圆心\((0,a)\)，\((a,0)\)，\(d=\sqrt{2}|a|\)，半径3、1．外离\(\Longleftrightarrow\sqrt{2}|a|>4\)，即\(|a|>2\sqrt{2}\)．\(-4\)满足，3满足（\(3>2\sqrt{2}\approx 2.83\)），\(-2\)与\(2\sqrt{2}\)均不满足．故选：AD．}
\end{ansblock}

\tihao{11}\zu{两圆公共弦与根轴×填空} \tieside{简单(知识点二)}若圆\(x^{2}+y^{2}=4\)与圆\(x^{2}+y^{2}+2ay-6=0\)（\(a>0\)）的公共弦长为\(2\sqrt{3}\)，则\(a=\)\kongda{1}.

\begin{ansblock}[2章-练-课时09-T11]
% ans:2章-练-课时09-T11
\ansitem{11}{1}
\ansnote{详解}{两方程相减得公共弦所在直线：\(2ay-2=0\)，即\(y=1/a\)，到圆心\(O\)的距离\(d=1/a\)．由半弦长平方\(3=4-1/a^{2}\)，得\(a^{2}=1\)，又\(a>0\)，故\(a=1\)．}
\end{ansblock}

\tihao{12}\zu{圆与圆的位置关系×填空} \tieside{简单(知识点三)}已知圆\(C\)：\((x-3)^{2}+(y-\sqrt{7})^{2}=1\)和两点\(A(-m,0)\)，\(B(m,0)\)（\(m>0\)），若圆\(C\)上存在点\(P\)，使得\(\angle APB=90^\circ\)，则\(m\)的最大值为\kongda{5}．

\begin{ansblock}[2章-练-课时09-T12]
% ans:2章-练-课时09-T12
\ansitem{12}{5}
\ansnote{详解}{\(\angle APB=90^\circ\Longleftrightarrow P\)在以\(AB\)为直径的圆\(x^{2}+y^{2}=m^{2}\)上．存在这样的\(P\Longleftrightarrow\)两圆有公共点：\(|m-1|\leqslant|OC|=\sqrt{9+7}=4\leqslant m+1\)．由\(4\geqslant m-1\)得\(m\leqslant 5\)（另一侧\(4\leqslant m+1\)给\(m\geqslant 3\)），故\(m\)的最大值为5．}
\end{ansblock}

\tihao{13}\zu{圆与圆的位置关系×填空} \tieside{中档(知识点一)}已知圆\(C_1\)的标准方程是\((x-4)^{2}+(y-4)^{2}=25\)，圆\(C_2\)：\(x^{2}+y^{2}-4x+my+3=0\)关于直线\(x+\sqrt{3}y+1=0\)对称，则圆\(C_1\)与圆\(C_2\)的位置关系为\kongda{相交}．

\begin{ansblock}[2章-练-课时09-T13]
% ans:2章-练-课时09-T13
\ansitem{13}{相交}
\ansnote{详解}{\(C_2\)圆心\((2,-m/2)\)在对称轴上：\(2-(\sqrt{3}/2)m+1=0\)，得\(m=2\sqrt{3}\)．\(r_2^{2}=4+3-3=4\)，\(r_2=2\)，圆心\(C_2(2,-\sqrt{3})\)；\(C_1\)圆心\((4,4)\)，\(r_1=5\)．\par\noindent \(d^{2}=4+(4+\sqrt{3})^{2}=23+8\sqrt{3}\approx 36.9\)，而\((r_1-r_2)^{2}=9\)，\((r_1+r_2)^{2}=49\)，\(9<23+8\sqrt{3}<49\)，即\(3<d<7\)，两圆相交．}
\end{ansblock}

\tihao{14}\zu{圆的切线方程×填空} \tieside{中档(知识点二)}已知圆\(C\)：\(x^{2}+y^{2}-2x-2y-2=0\)，直线\(l\)：\(x+2y+2=0\)，\(M\)为直线\(l\)上的动点，过点\(M\)作圆\(C\)的切线\(MA\)，\(MB\)，切点为\(A\)，\(B\)，当四边形\(MACB\)的面积取最小值时，直线\(AB\)的方程为\kongda{x+2y+1=0}．

\begin{ansblock}[2章-练-课时09-T14]
% ans:2章-练-课时09-T14
\ansitem{14}{x+2y+1=0}
\ansnote{详解}{圆即\((x-1)^{2}+(y-1)^{2}=4\)，圆心\((1,1)\)，\(r=2\)．\(S=2\cdot(1/2)r\cdot|MA|=2\sqrt{|MC|^{2}-4}\)，最小\(\Longleftrightarrow|MC|\)最小，即\(C\)到\(l\)的距离\(|1+2+2|/\sqrt{5}=\sqrt{5}\)．\par\noindent 此时\(M\)为垂足：\(CM\)方向\((1,2)\)，直线\(CM\)：\(y-1=2(x-1)\)，与\(l\)联立得\(M(0,-1)\)．以\(CM\)为直径的圆\(x^{2}+y^{2}-x-1=0\)与圆\(C\)相减得\(x+2y+1=0\)，即切点弦\(AB\)的方程．}
\end{ansblock}

% ============================================================
% 思维探索（题15~16＝T15~T16：解答2，居末档照库序）
% ============================================================
\dangbadge{思}{维}{探}{索}

\tihao{15}\zu{两圆公共弦与根轴×解答} \tieside{简单(知识点二)}某同学发现了一个现象：在求圆的公共弦\(AB\)（即两个圆相交时，两个交点的连线）所在直线的方程恰好与两个圆的方程相减消掉二次项\(x^{2}\)，\(y^{2}\)后所得的方程一样．由此，他提出了一个猜想：对于两个圆\(C_1\)：\(x^{2}+y^{2}+D_1x+E_1y+F_1=0\)与\(C_2\)：\(x^{2}+y^{2}+D_2x+E_2y+F_2=0\)，直线\((D_1-D_2)x+(E_1-E_2)y+(F_1-F_2)=0\)就是两个圆的公共弦所在直线的方程．你认为他的猜想对吗？请说明理由．
\liubai[16mm]{解答书写区}

\begin{ansblock}[2章-练-课时09-T15]
% ans:2章-练-课时09-T15
\ansitem{15}{猜想正确，理由见解析}
\ansnote{详解}{设交点\(A(x_1,y_1)\)，\(B(x_2,y_2)\)，则\(A\)、\(B\)坐标同时满足两圆方程：\(x_i^{2}+y_i^{2}+D_1x_i+E_1y_i+F_1=0\)与\(x_i^{2}+y_i^{2}+D_2x_i+E_2y_i+F_2=0\)（\(i=1,2\)）．\par\noindent 两式相减：\((D_1-D_2)x_i+(E_1-E_2)y_i+(F_1-F_2)=0\)，即\(A\)、\(B\)都在直线\((D_1-D_2)x+(E_1-E_2)y+(F_1-F_2)=0\)上．又\(D_1\neq D_2\)或\(E_1\neq E_2\)（否则两圆方程相同），该直线系数不全为零，而过两个不同交点的直线有且仅有一条，故它就是公共弦\(AB\)所在直线的方程，猜想正确．}
\end{ansblock}

\tihao{16}\zu{圆与圆的位置关系×解答} \tieside{简单(知识点二)}求过点\(A(1,-1)\)且与圆\(C\)：\(x^{2}+y^{2}=100\)相切于点\(B(8,6)\)的圆的方程．
\liubai[16mm]{解答书写区}

\begin{ansblock}[2章-练-课时09-T16]
% ans:2章-练-课时09-T16
\ansitem{16}{(x−4)²+(y−3)²=25}
\ansnote{详解}{设圆心\(D(a,b)\)．切点\(B\)在两圆连心线上（连心线过切点且与公切线垂直）：\(k_{CB}=k_{DB}\)，即\(b/a=6/8=3/4\)，\(b=3a/4\)．又\(|DA|=|DB|\)：\((a-1)^{2}+(b+1)^{2}=(a-8)^{2}+(b-6)^{2}\)，展开得\(14a+14b=98\)，即\(a+b=7\)．联立\(b=3a/4\)得\(a=4\)，\(b=3\)，\(r=|DB|=\sqrt{16+9}=5\)．\par\noindent 故所求圆的方程为\((x-4)^{2}+(y-3)^{2}=25\)．（验：\(|DA|=\sqrt{9+16}=5\)；\(|DC|=10-5=5\)，两圆内切于\(B\)．）}
\end{ansblock}

% ---- 尾块（\tailfill 置 \end{multicols} 前；无参＝内容尾随框，每件恰 1 处） ----
\tailfill

\end{multicols}
\end{document}
